
dans les deux cas envisag\'es. Comme \(\Pic^\tau_{Y/k}\) est propre, \(Y\) \'etant normal (FGA 236-2.13), et comme \(N\) est un sous-groupe, n\'ecessairement ferm\'e (SGA 3 VI\(_B\) 1.9.2), pour que \(N\) soit fini, il suffit qu'il ne contienne pas le sous-groupe de type \(\mu_\ell(=\mathbb Z/\ell\mathbb Z)\) o\`u \((\ell,\operatorname{car}(k))=1\).

Consid\'erons la th\'eorie de Kummer, c'est-\`a-dire, le diagramme commutatif suivant, provenant de la suite exacte
\[
0\longrightarrow\mu_\ell\longrightarrow\mathbb G_m
\xrightarrow{x\mapsto x^\ell}\mathbb G_m\longrightarrow0
\]
de faisceaux \'etales sur \(Y\) (resp. sur \(X\))~:
\[
\begin{tikzcd}[column sep=1.6em,row sep=1.2em]
& & H^1(Y,\mu_\ell)\arrow[r]\arrow[d]&H^1(Y,\mathbb G_m)\arrow[r,"\ell"]\arrow[d]&H^1(Y,\mathbb G_m)\arrow[d]\\
H^0(X,\mathbb G_m)\arrow[r,"\ell"]&H^0(X,\mathbb G_m)\arrow[r]&H^1(X,\mu_\ell)\arrow[r]&H^1(X,\mathbb G_m)\arrow[r,"\ell"]&H^1(X,\mathbb G_m).
\end{tikzcd}
\]
Par le ``th\'eor\`eme 90'' de Hilbert, on a \(H^1(Y,\mathbb G_m)=\Pic_{Y/k}(k)\) (resp. \(H^1(X,\mathbb G_m)=\Pic_{X/k}(k)\)). Comme \(X\) est propre et int\`egre, on a \(H^0(X,\mathbb G_m)=k^*\)\,; donc \(H^1(X,\mu_\ell)\to H^1(X,\mathbb G_m)\) est injective.

Il suffit donc de voir que \(H^1(Y,\mathbb Z/\ell\mathbb Z)\longrightarrow H^1(X,\mathbb Z/\ell\mathbb Z)\) est injectif, ce qui r\'esulte du fait que \(\pi_1(X,\xi)\to\pi_1(Y,\xi)\) est surjectif, ou encore (SGA 1 V 6.9) que tout rev\^etement \'etale connexe \(\widetilde Y\) de \(Y\) de groupe \(\mathbb Z/\ell\mathbb Z\) a une restriction \(\widetilde X=X\times_Y\widetilde Y\) qui est aussi connexe, ce qui est une cons\'equence bien connue du th\'eor\`eme de Bertini et du th\'eor\`eme de connexion de Zariski (SGA 1 X 2.10).

\textbf{Remarque 3.12.} Conservons les conditions de 3.11. Alors :
\begin{enumerate}[label=(\roman*)]
\item Le noyau \(N\) est un groupe fini \uline{unipotent} (SGA 3 XVII 1.1), puisqu'il ne contient pas le sous-groupe de type multiplicatif \(\mu_n\) pour aucun \(n\) (SGA 3 XVII 4.6.1 vi), ce qu'on prouve de m\^eme, utilisant au
